% Options for packages loaded elsewhere
\PassOptionsToPackage{unicode}{hyperref}
\PassOptionsToPackage{hyphens}{url}
\documentclass[
]{article}
\usepackage{xcolor}
\usepackage[margin=1in]{geometry}
\usepackage{amsmath,amssymb}
\setcounter{secnumdepth}{-\maxdimen} % remove section numbering
\usepackage{iftex}
\ifPDFTeX
  \usepackage[T1]{fontenc}
  \usepackage[utf8]{inputenc}
  \usepackage{textcomp} % provide euro and other symbols
\else % if luatex or xetex
  \usepackage{unicode-math} % this also loads fontspec
  \defaultfontfeatures{Scale=MatchLowercase}
  \defaultfontfeatures[\rmfamily]{Ligatures=TeX,Scale=1}
\fi
\usepackage{lmodern}
\ifPDFTeX\else
  % xetex/luatex font selection
\fi
% Use upquote if available, for straight quotes in verbatim environments
\IfFileExists{upquote.sty}{\usepackage{upquote}}{}
\IfFileExists{microtype.sty}{% use microtype if available
  \usepackage[]{microtype}
  \UseMicrotypeSet[protrusion]{basicmath} % disable protrusion for tt fonts
}{}
\makeatletter
\@ifundefined{KOMAClassName}{% if non-KOMA class
  \IfFileExists{parskip.sty}{%
    \usepackage{parskip}
  }{% else
    \setlength{\parindent}{0pt}
    \setlength{\parskip}{6pt plus 2pt minus 1pt}}
}{% if KOMA class
  \KOMAoptions{parskip=half}}
\makeatother
\usepackage{longtable,booktabs,array}
\usepackage{calc} % for calculating minipage widths
% Correct order of tables after \paragraph or \subparagraph
\usepackage{etoolbox}
\makeatletter
\patchcmd\longtable{\par}{\if@noskipsec\mbox{}\fi\par}{}{}
\makeatother
% Allow footnotes in longtable head/foot
\IfFileExists{footnotehyper.sty}{\usepackage{footnotehyper}}{\usepackage{footnote}}
\makesavenoteenv{longtable}
\usepackage{graphicx}
\makeatletter
\newsavebox\pandoc@box
\newcommand*\pandocbounded[1]{% scales image to fit in text height/width
  \sbox\pandoc@box{#1}%
  \Gscale@div\@tempa{\textheight}{\dimexpr\ht\pandoc@box+\dp\pandoc@box\relax}%
  \Gscale@div\@tempb{\linewidth}{\wd\pandoc@box}%
  \ifdim\@tempb\p@<\@tempa\p@\let\@tempa\@tempb\fi% select the smaller of both
  \ifdim\@tempa\p@<\p@\scalebox{\@tempa}{\usebox\pandoc@box}%
  \else\usebox{\pandoc@box}%
  \fi%
}
% Set default figure placement to htbp
\def\fps@figure{htbp}
\makeatother
\setlength{\emergencystretch}{3em} % prevent overfull lines
\providecommand{\tightlist}{%
  \setlength{\itemsep}{0pt}\setlength{\parskip}{0pt}}
\usepackage{bookmark}
\IfFileExists{xurl.sty}{\usepackage{xurl}}{} % add URL line breaks if available
\urlstyle{same}
\hypersetup{
  hidelinks,
  pdfcreator={LaTeX via pandoc}}

\author{}
\date{\vspace{-2.5em}}

\begin{document}

\subsection{\#
DataCollection\_analysis\_fandom}\label{datacollection_analysis_fandom}

title: ``Fandom Association with loneliness analysis'' The group work on
the analysis of correlations between interactions in fandom and the
level of loneliness. Here we import the excel dataset into the RStudio
software first. ---

\texttt{\{r\}\ library(dplyr)\ library(readr)\ library(ggplot2)\ d\_loneliness\_association\ \textless{}-\ read.csv(\textquotesingle{}Loneliness\_association.csv\textquotesingle{})}

\begin{longtable}[]{@{}l@{}}
\toprule\noalign{}
\endhead
\bottomrule\noalign{}
\endlastfoot
First, we exclude all the answers that say they don't follow any
celebrities in real life. \\
\end{longtable}

\texttt{\{r\}\ d\_loneliness\_association\ \textless{}-\ d\_loneliness\_association\ \%\textgreater{}\%\ \ \ filter(\ \ \ \ \ Follow\_celeb\ ==\ \textquotesingle{}Yes.\textquotesingle{}\ \ \ )}
--- Then we could then convert all our ordinal answers into numerical
forms that allow us to add up and summarise the scores. Here, we sort
out the series of questions that test the loneliness, L\_C1 - L\_C8,
with two reverse questions. Here the higher the score is, the higher the
loneliness level would be. ---

\texttt{\{r\}\ normal\_questions\ \textless{}-\ c(\textquotesingle{}L\_C1\textquotesingle{},\ \textquotesingle{}L\_C2\textquotesingle{},\ \textquotesingle{}L\_C4\textquotesingle{},\ \textquotesingle{}L\_C5\textquotesingle{},\ \textquotesingle{}L\_C7\textquotesingle{},\ \textquotesingle{}L\_C8\textquotesingle{})\ reverse\_questions\ \textless{}-\ c(\textquotesingle{}L\_C3\textquotesingle{},\ \textquotesingle{}L\_C6\textquotesingle{})\ d\_loneliness\_association\ \textless{}-\ d\_loneliness\_association\%\textgreater{}\%\ \ \ mutate(\ \ \ \ \ across(\ \ \ \ \ \ \ all\_of(normal\_questions),\ \ \ \ \ \ \ \textasciitilde{}case\_when(\ \ \ \ \ \ \ \ \ .\ ==\ \textquotesingle{}Never\textquotesingle{}\ \textasciitilde{}\ 1,\ \ \ \ \ \ \ \ \ .\ ==\ \textquotesingle{}Rarely\textquotesingle{}\ \textasciitilde{}\ 2,\ \ \ \ \ \ \ \ \ .\ ==\ \textquotesingle{}Sometimes\textquotesingle{}\ \textasciitilde{}\ 3,\ \ \ \ \ \ \ \ \ .\ ==\ \textquotesingle{}Often\textquotesingle{}\ \textasciitilde{}\ 4,\ \ \ \ \ \ \ \ \ TRUE\ \textasciitilde{}\ NA\_real\_\ \ \ \ \ \ \ ),\ \ \ \ \ \ \ .names\ =\ "\{.col\}\_num"\ \ \ \ \ ),\ \ \ \ \ across(\ \ \ \ \ \ \ all\_of(reverse\_questions),\ \ \ \ \ \ \ \textasciitilde{}case\_when(\ \ \ \ \ \ \ \ \ .\ ==\ \textquotesingle{}Never\textquotesingle{}\ \textasciitilde{}\ 4,\ \ \ \ \ \ \ \ \ .\ ==\ \textquotesingle{}Rarely\textquotesingle{}\ \textasciitilde{}\ 3,\ \ \ \ \ \ \ \ \ .\ ==\ \textquotesingle{}Sometimes\textquotesingle{}\ \textasciitilde{}\ 2,\ \ \ \ \ \ \ \ \ .\ ==\ \textquotesingle{}Often\textquotesingle{}\ \textasciitilde{}\ 1,\ \ \ \ \ \ \ \ \ TRUE\ \textasciitilde{}\ NA\_real\_\ \ \ \ \ \ \ ),\ \ \ \ \ \ \ .names\ =\ "\{.col\}\_num"\ \ \ \ \ )\ \ \ )}
--- Here we then clean out the questions on the participation with
fandom, marked F\_C1 - F\_C7, with no reverse questions. Here, the
higher the scores are, the more participated you are in the fandoms. ---

\texttt{\{r\}\ questions\ \textless{}-\ c(\textquotesingle{}F\_C1\textquotesingle{},\ \textquotesingle{}F\_C2\textquotesingle{},\ \textquotesingle{}F\_C3\textquotesingle{},\ \textquotesingle{}F\_C4\textquotesingle{},\ \textquotesingle{}F\_C5\textquotesingle{},\ \textquotesingle{}F\_C6\textquotesingle{},\ \textquotesingle{}F\_C7\textquotesingle{})\ d\_loneliness\_association\ \textless{}-\ d\_loneliness\_association\%\textgreater{}\%\ \ \ mutate(\ \ \ \ \ across(\ \ \ \ \ \ \ all\_of(questions),\ \ \ \ \ \ \ \textasciitilde{}case\_when(\ \ \ \ \ \ \ \ \ .==\ "Strongly\ disagree"\ \textasciitilde{}\ 1,\ \ \ \ \ \ \ \ \ .==\ "Disagree"\ \textasciitilde{}\ 2,\ \ \ \ \ \ \ \ \ .==\ "Neutral"\ \textasciitilde{}\ 3,\ \ \ \ \ \ \ \ \ .==\ "Agree"\ \textasciitilde{}\ 4,\ \ \ \ \ \ \ \ \ .==\ "Strongly\ agree"\ \textasciitilde{}\ 5,\ \ \ \ \ \ \ \ \ TRUE\ \textasciitilde{}\ NA\_real\_\ \ \ \ \ \ \ ),\ \ \ \ \ \ \ .names\ =\ "\{.col\}\_num"\ \ \ \ \ )\ \ \ )}
--- Then we can sum up all the columns together to get two overall
scores. ---
\texttt{\{r\}\ d\_loneliness\_association\ \textless{}-\ d\_loneliness\_association\%\textgreater{}\%\ \ \ mutate(\ \ \ \ \ frequency\_score\ =\ rowSums(\ \ \ \ \ \ \ across(all\_of(paste0(c(normal\_questions,\ reverse\_questions),\ "\_num"))),\ \ \ \ \ \ \ na.rm\ =\ TRUE\ \ \ ),\ \ \ agreement\_score\ =\ rowSums(\ \ \ \ \ \ \ across(all\_of(paste0(c(questions),\ "\_num"))),\ \ \ \ \ \ \ na.rm\ =\ TRUE\ \ \ \ \ )\ \ \ )}
--- After the cleaning, we can see the distribution of gender and age
respectively, starting with gender first, then age groups. The gender is
done only on biological gender instead of the identity. We could also
calculate the mean of both scores to see the average participation level
and the average degree of loneliness. ---

```\{r\} library(psych)

Frequency\_questions\_num \textless- c(`L\_C1\_num', `L\_C2\_num',
`L\_C3\_num', `L\_C4\_num', `L\_C5\_num',`L\_C6\_num', `L\_C7\_num',
`L\_C8\_num') Agreement\_questions\_num \textless- c(`F\_C1\_num',
`F\_C2\_num', `F\_C3\_num', `F\_C4\_num', `F\_C5\_num', `F\_C6\_num',
`F\_C7\_num') psych::alpha(d\_loneliness\_association{[},
Frequency\_questions\_num{]})

psych::alpha(d\_loneliness\_association{[},
Agreement\_questions\_num{]})

\begin{verbatim}

```{r}
d_loneliness_association %>%
  group_by(
    Gender
  ) %>%summarise(
    n = n(),
    M1 = mean(frequency_score),
    M2 = mean(agreement_score)
    )
\end{verbatim}

\begin{longtable}[]{@{}l@{}}
\toprule\noalign{}
\endhead
\bottomrule\noalign{}
\endlastfoot
Here above we can see that there is only 1 male participant. \\
\end{longtable}

\texttt{\{r\}\ d\_loneliness\_association\ \%\textgreater{}\%\ \ \ group\_by(\ \ \ \ \ Age\_group\ \ \ )\ \%\textgreater{}\%summarise(\ \ \ \ \ n\ =\ n(),\ \ \ \ \ M1\ =\ mean(frequency\_score),\ \ \ \ \ M2\ =\ mean(agreement\_score),\ \ \ \ \ SD1\ =\ sd(frequency\_score),\ \ \ \ \ SD2\ =\ sd(agreement\_score),\ \ \ \ \ MED1\ =\ median(frequency\_score),\ \ \ \ \ MED2\ =\ median(agreement\_score)\ \ \ \ \ )}

\texttt{\{r\}\ d\_loneliness\_association\ \%\textgreater{}\%\ \ \ ggplot(\ \ \ \ \ \ \ aes(x\ =\ frequency\_score,\ y\ =\ agreement\_score)\ \ \ )\ +\ \ \ \ geom\_jitter()}

\texttt{\{r\}\ d\_loneliness\_association\ \%\textgreater{}\%\ \ \ cor.test(\ \ \ \ \ \textasciitilde{}\ agreement\_score\ +\ frequency\_score,\ \ \ \ \ data\ =\ .\ \ \ )}

\end{document}
